% Separate proved research note.  This does not change the frozen C.90 source.
% No claim number has been assigned to this note.

\paragraph{Cocompact full-group refinement (proved).}
Retain the actual hypotheses of C.81 and the compact quotient
$N=O/Q$ of C.90.  Thus $O=\operatorname{int}\overline D$ is
the original open disk, $Q$ is closed, discrete, free and proper,
and an actual subgroup of the boundary-fixed centralizer
$\mathscr C_Q$ is transitive on $O$.
Fix $x\in O$, put $p=\pi x$, and write
$\mathscr F=\mathscr C_Q(x)$ and
$\mathscr H_p=\operatorname{Homeo}(N\mathbin{\mathrm{rel}}\{p\})^0$.
Give these groups their stated compact-open map-and-inverse
topologies.  There are exact topological group isomorphisms
\[
 \begin{aligned}
 D:\mathscr F&\longrightarrow\mathscr H_p,
       &D(h)(\pi z)&=\pi(hz),\\
 L:\mathscr H_p&\longrightarrow\mathscr F,
       &L(b)|_{\partial\overline D}&=\id.
 \end{aligned}                                                    \tag{CG.1}
\]
The interior restriction of $L(b)$ is exactly the unique
lift of $b\pi$ that fixes $x$.
In particular this identifies the entire actual point stabilizer,
not only its identity component.  The proof establishes that every
one of its descents lies in the indicated pointed component.
The evaluation map admits a continuous section
$s:O\to\mathscr C_Q$ with $s(y)x=y$ and $s(x)=1$.
Consequently the mutually inverse homeomorphisms
\[
 \begin{aligned}
 \Phi:O\times\mathscr F&\longrightarrow\mathscr C_Q,
       &\Phi(y,k)&=s(y)k,\\
 \Psi:\mathscr C_Q&\longrightarrow O\times\mathscr F,
       &\Psi(h)&=(hx,s(hx)^{-1}h)
 \end{aligned}                                                    \tag{CG.2}
\]
give an identity-fixed contraction of the whole original
$\mathscr C_Q$.  All maps in the contraction fix the original
circle and commute with every individual original $q\in Q$.
The assertion includes the torus, Klein bottle, and all the
other compact quotient surfaces permitted by these data.

\begin{proof}
\emph{A specified plane coordinate used only for homotopies and cells.}
Choose a disk coordinate
$\theta:\overline D\to\{u\in\mathbb C:|u|\leq1\}$ and put
$a=\theta(x)$.  Define
\[
 \begin{aligned}
 m_a(u)&=\frac{u-a}{1-\overline a u},
       &m_a^{-1}(v)&=\frac{v+a}{1+\overline a v},\\
 \psi(z)&=\frac{m_a(\theta(z))}{1-|m_a(\theta(z))|},
       &\psi^{-1}(w)&=
          \theta^{-1}m_a^{-1}\!\left(\frac{w}{1+|w|}\right).
 \end{aligned}                                                    \tag{CG.3}
\]
These displayed inverse formulas give a homeomorphism
$\psi:O\to\mathbb R^2$ with $\psi(x)=0$.
They are a coordinate on the specified interior; all operator
equations below continue to use the original $Q$ and the
original quotient covering.  In particular the coordinate
does not replace the original boundary topology of the group.

\emph{Every actual pointed descent is pointed homotopic to the identity.}
For $h\in\mathscr F$, the formulas
\[
             b(\pi z)=\pi(hz),\qquad
             b^{-1}(\pi z)=\pi(h^{-1}z)                         \tag{CG.4}
\]
are well-defined, mutually inverse continuous maps, because
$h$ and $h^{-1}$ commute with every $q$.  Both fix $p$.
For clarity, the induced fundamental-group comparison retains
the original labels.  A loop $\gamma$ at $p$ has a unique lift
$\widetilde\gamma$ starting at $x$, ending at some $qx$.
Its image $h\widetilde\gamma$ lifts $b\gamma$, starts at $x$,
and ends at $h(qx)=qh(x)=qx$.  Lifts with the same endpoints
are homotopic rel endpoints in the simply connected $O$.
Projecting such a homotopy shows
$[b\gamma]=[\gamma]$ in the actual $\pi_1(N,p)$.

Here is a direct based homotopy, so no classification of
pointed mapping classes is being presumed from that computation.
Take a finite triangulation of $N$ in which $p$ is a vertex,
as in C.90, and lift it through $\pi$.
Every lifted simplex maps homeomorphically onto its quotient
simplex.  A deck map preserving such a lift fixes a point
of it by injectivity of that restriction and hence is the
identity.  Choose one lift of every quotient simplex,
with the chosen lift of $p$ equal to $x$.
We construct a $Q$-equivariant homotopy
\[
 W:O\times[0,1]\longrightarrow O,
 \qquad W_0=h|_O,\quad W_1=\id_O,\quad W_t(qx)=qx.            \tag{CG.5}
\]
On a chosen vertex $v$, use the path
$\psi^{-1}((1-t)\psi(hv)+t\psi(v))$.
For $v=x$ this is constantly $x$.
Define the paths on every translate by the literal formula
$W_t(qv)=qW_t(v)$.

For a chosen lifted edge $e$, the paths on its two endpoints,
together with $h|_e$ at time zero and $\id|_e$ at time one,
give a continuous map on the boundary of the square
$e\times[0,1]$.  This extends across the square.
An explicit extension for any boundary map $f:S^j\to O$
is obtained in cone coordinates by
\[
       (u,r)\longmapsto\psi^{-1}(r\psi(f(u))),
       \qquad 0\leq r\leq1,                                  \tag{CG.6}
\]
with the cone vertex sent to $x$; a fixed ball-to-prism
homeomorphism supplies the square coordinates here.
Continuity at the cone vertex follows because $f(S^j)$
is compact and $\psi(f(S^j))$ is bounded.
Transport this extension to every translate of $e$ by $q$.
It agrees with the already defined vertex paths.

For a chosen lifted triangle, the two end faces and the
three edge homotopies define a continuous map on the
boundary of its triangular prism.  Apply (CG.6) with
$j=2$ and a fixed ball-to-prism homeomorphism, and transport
the extension to all its $Q$-translates.
The boundary prescriptions agree on every common face.
The lifted triangulation is locally finite, so the closed-cell
pasting lemma on $O\times[0,1]$ proves continuity of $W$.
Each equivariance equation holds on every cell and hence
on all of $O$.  The vertex prescription proves the final
fixed-orbit equation in (CG.5).

It follows that
\[
             B_t(\pi z)=\pi W_t(z)                            \tag{CG.7}
\]
is a well-defined continuous homotopy from $b$ to $\id_N$
fixing $p$ at every time.  Continuity follows, for example,
by using local sections of the original covering.
This homotopy is proper: its domain $N\times[0,1]$ is
compact and the inverse image of a compact subset of
the Hausdorff surface $N$ is closed in that domain.
The surface $N$ has empty boundary, so the boundary-pair
condition in Epstein's theorem is satisfied as well.

Apply exactly \cite[Theorem~6.3, p.~99; proof, pp.~100--101]
{C_Epstein1966}.  This is the theorem for a proper homotopy
of the boundary pair; its last sentence retains a basepoint
fixed by the homotopy.  Its proof explicitly treats a
non-PL homotopy on p.~101 by a basepoint-preserving isotopy
and PL approximation, using the Appendix; the stated
topological conclusion therefore applies to our actual $b$.
We obtain an isotopy from $b$ to $\id_N$ fixing $p$.
On the compact $N$, the isotopy is continuous in the
compact-open topology of maps, and its inverse family is
continuous as well: the level map on $N\times[0,1]$ is
a continuous bijection of a compact space onto a Hausdorff
space.  Thus $b\in\mathscr H_p$ in the stated group topology.
This proves the asserted component membership for every
$h\in\mathscr F$, without first restricting $h$ to a component.

\emph{The exact pointed compact-surface contraction and all exceptions.}
As in C.90, $N$ is a closed connected PL surface with
universal cover $O$.  The covering isomorphism on second
homotopy groups gives $\pi_2(N)=0$, whereas $S^2$ and
$\mathbb{RP}^2$ have second homotopy group $\mathbb Z$.
These two closed-surface cases are therefore excluded
by the actual covering.
For the pair $(N,\{p\})$, Yagasaki's
\cite[Lemma~2.1(i), p.~2; proof, p.~3]{C_Yagasaki2000}
proves contractibility of $\mathscr H_p$.
The exceptional disk, annulus, and M\"obius-band pairs
have boundary and cannot be our $N$.
The exceptional torus and Klein-bottle pairs have empty
fixed set, whereas our fixed set is $\{p\}$.
The remaining listed closed exceptions belong to the
already excluded sphere and projective plane.
No orientability assumption occurs in this application.
Likewise Epstein's Theorem~6.3, used above with the
proper based homotopy (CG.7), includes the closed
nonorientable surfaces.  We have not used the distinct
nonproper-homotopy statement of his Theorem~6.4(b)
or its open or closed M\"obius-band qualification.

Write $C_t$ for a contraction of $\mathscr H_p$,
with $C_0(b)=b$ and constant endpoint $C_1(b)=c$.
The exact normalization already used in C.90 is
\[
 J_r(b)=C_{1-r}(b)C_{1-r}(1)^{-1},
 \qquad J_0(b)=1,\quad J_1(b)=b,\quad J_r(1)=1.                \tag{CG.8}
\]
Every value and every factor fixes $p$.

\emph{The inverse lift and the two actual topologies.}
For $b\in\mathscr H_p$, lift the homotopy
$J_r(b)\pi:O\to N$ from the identity lift at $r=0$.
Let $H_r(b)$ be the resulting interior family.
Lift $J_r(b)^{-1}\pi$ from the identity to obtain $K_r(b)$.
The two compositions $H_r(b)K_r(b)$ and $K_r(b)H_r(b)$
lift the constant homotopy $\pi$ and start at the identity;
uniqueness makes both compositions the identity.
For every literal $q\in Q$, the two maps
$H_r(b)q$ and $qH_r(b)$ lift the same homotopy and
start at the same map $q$, so they agree.
The track of $x$ lies in the discrete fiber over $p$
and starts at $x$, hence remains $x$.

C.88 applies to this compact parameter interval with
$K=N$: the required tail is empty.
It supplies the continuous original-circle identity
extensions of $H_r(b)$ and $K_r(b)$, which remain inverse.
Define $L(b)$ as the extension of $H_1(b)$.
It is therefore an actual member of $\mathscr F$.
It descends to $b$.  Any two lifts of $b\pi$ that take
$x$ to $x$ agree on the connected $O$; their boundary
extensions agree too.  This proves
\[
              D L(b)=b,\qquad L D(h)=h.                       \tag{CG.9}
\]
The same uniqueness gives $L(bc)=L(b)L(c)$, so both maps
in (CG.1) are group maps with the displayed inverse.

The descent topology has a direct metric estimate.
Use the actual metrics $d,d_Q,\rho$ of C.81 and set
\[
 \begin{aligned}
 d_{\mathscr C}(h,k)
   &=\sup_{z\in\overline D}d(hz,kz)
      +\sup_{z\in\overline D}d(h^{-1}z,k^{-1}z),\\
 d_{\mathscr H}(b,c)
   &=\sup_{u\in N}\rho(bu,cu)
      +\sup_{u\in N}\rho(b^{-1}u,c^{-1}u).
 \end{aligned}
\]
They induce the asserted map-and-inverse topologies
because their underlying spaces are compact.
The quotient inequality and the exact C.81 isometry give
\[
                 d_{\mathscr H}(Dh,Dk)
                       \leq d_{\mathscr C}(h,k).              \tag{CG.10}
\]
Thus $D$ is continuous in the original group topology.

For the inverse, take any sequence $b_n\to b$ in
$\mathscr H_p$ and use the compact metric parameter space
$P=\{0\}\cup\{1/n:n\geq1\}$ with $b_0=b$.
The family $J_r(b_u)$ and its inverse are continuous
on $[0,1]\times P\times N$.
Lift them jointly from the identity through the original
covering, exactly as above.  The resulting interior
families have every original commutation equation and
fix $x$.  C.88, now with compact parameter
$[0,1]\times P$ and again $K=N$, supplies joint
continuity of their identity extensions and inverses
on the original closed disk.  At $r=1$ this gives
$L(b_n)\to L(b)$ in $d_{\mathscr C}$.
Both groups in (CG.1) are metrizable subspaces of
the indicated compact-domain homeomorphism groups,
so this sequential argument proves continuity of $L$.
It proves the entire asserted topological group
isomorphism, rather than only a set-theoretic lifting rule.

In particular the single formula
\[
 A_t(k)=L\bigl(J_{1-t}(Dk)\bigr),
 \quad A_0(k)=k,\quad A_1(k)=1,\quad A_t(1)=1                \tag{CG.11}
\]
is a continuous identity-fixed contraction of the whole
point stabilizer $\mathscr F$ in its original topology.

\emph{A global section with explicit local formulas.}
C.81 supplies actual continuous local sections of
$h\mapsto hx$ over an open cover $\{U_i\}$ of $O$.
Write a chosen one on $U_i$ as
$\sigma_i:U_i\to\mathscr C_Q$, so that
$\sigma_i(y)x=y$; these maps already have the original
map-and-inverse topology by C.81.11.
Start with the standard locally finite unit-grid
triangulation of $\mathbb R^2$, with $0$ a vertex.
Apply exactly the $k=0$ subdivision lemma used in C.82,
\cite[Lemma~1, p.~458]{C_Haver1975}, to this complex,
the actual map $\psi^{-1}$ into $O$, and the cover
$\{U_i\}$.  It gives a locally finite subdivision
whose closed vertex stars map into members of that
cover.  Each closed simplex is in the closed star
of any one of its vertices, so its image lies in
some $U_i$.  Pull the subdivided triangles back to
$O$ by $\psi^{-1}$.  The original vertex $0$ remains
a vertex under subdivision, and its image is $x$.

At every vertex $v$, choose once an actual
$s(v)\in\mathscr C_Q$ taking $x$ to $v$,
using the surjectivity in C.81.  Choose $s(x)=1$.
We extend this section over the edges and then over
the two-simplices.  Suppose its values on the boundary
of a closed simplex $\tau$ have been defined.
Choose $i(\tau)$ with $\tau\subset U_{i(\tau)}$ and put
\[
 a_\tau(u)=\sigma_{i(\tau)}(u)^{-1}s(u)\in\mathscr F
                         \qquad(u\in\partial\tau).            \tag{CG.12}
\]
In the affine coordinates of the plane simplex, let
$c_\tau$ be its barycenter.  Every noncentral point
has a unique expression $y=(1-r)c_\tau+ru$ with
$u\in\partial\tau$ and $0<r\leq1$; this notation
is transported to $O$ by $\psi^{-1}$.
Define the extension by the exact formulas
\[
 \begin{aligned}
 s(y)&=\sigma_{i(\tau)}(y)
                   A_{1-r}\bigl(a_\tau(u)\bigr)
           &&(0<r\leq1),\\
 s(c_\tau)&=\sigma_{i(\tau)}(c_\tau).
 \end{aligned}                                                    \tag{CG.13}
\]
The second factor always fixes $x$, so $s(y)x=y$.
At $r=1$ its value is exactly the already prescribed
boundary value.  At the center, continuity follows
from $A_1=1$ uniformly on the compact set
$a_\tau(\partial\tau)$: the restriction of the
continuous map $A$ to
$[0,1]\times a_\tau(\partial\tau)$ is uniformly
continuous in any compatible metric on its image.
This checks the possible polar-coordinate singularity
without imposing any uniform condition at infinity.
The edge extensions fit at their endpoints, and
the face extensions retain every edge value.
Local finiteness and the closed-cell pasting lemma
therefore give one global continuous
$s:O\to\mathscr C_Q$ in the actual group topology.
The original value $s(x)=1$ is retained.

Now the maps in (CG.2) are defined on the displayed
spaces: $s(hx)^{-1}h$ fixes $x$.
Their two composites are identities by multiplication
and $s(y)x=y$.  Evaluation, multiplication, inversion,
and $s$ are continuous in the original topology,
so these are mutually inverse homeomorphisms.
This also proves that the constructed section is a
section of the actual C.81 bundle, with its original
fiber and topology.

\emph{The contraction of the whole original group.}
In the specified coordinate (CG.3), put
\[
             R_t(y)=\psi^{-1}((1-t)\psi(y)).                   \tag{CG.14}
\]
This is a continuous contraction of $O$ to $x$,
fixing $x$ throughout.  For $h\in\mathscr C_Q$ let
$y=hx$ and $k=s(y)^{-1}h$.  The required global
contraction and its pointwise inverse map are
\[
 \begin{aligned}
 \mathcal C_t(h)&=s(R_t(y))A_t(k),\\
 \mathcal C_t(h)^{-1}&=A_t(k)^{-1}s(R_t(y))^{-1}.
 \end{aligned}                                                    \tag{CG.15}
\]
All factors belong to the original $\mathscr C_Q$,
so every value fixes its original circle pointwise
and commutes with each original $q$.
The formulas and (CG.2), (CG.11), (CG.14) prove
joint continuity in the map-and-inverse topology.
At time zero the first formula is $s(y)k=h$;
at time one it is $s(x)\cdot1=1$.
For $h=1$, we have $y=x$, $k=1$, so
$\mathcal C_t(1)=1$ for every $t$.
This proves the claimed contraction of the entire group.
\end{proof}
